% ==============================================================================
% Fichier     : main_exercices.tex
% Titre       : Résumé académique -- Exercices Fondamentaux
% Institution : Université de Yaoundé I -- ENSP
% ==============================================================================

\documentclass[12pt, a4paper]{report}

\usepackage{style}
\usepackage[utf8]{inputenc}
\usepackage[T1]{fontenc}
\usepackage[french]{babel}
\usepackage{multicol}
\sloppy

% ==============================================================================
% VARIABLES
% ==============================================================================
\newcommand{\Universite}{Université de Yaoundé I}
\newcommand{\Etablissement}{École Nationale Supérieure Polytechnique}
\newcommand{\Departement}{Département de Génie Informatique}
\newcommand{\AnneeAcademique}{2025 -- 2026}
\newcommand{\LogoEtab}{logo_polytech.png}
\newcommand{\CodeUE}{GES4062}
\newcommand{\TitreRapport}{TECHNIQUES D'INVESTIGATION NUMERIQUE AVANCEE}
\newcommand{\ThemeProjet}{\LARGE TEST DES COMMANDES SUR WINDOWS ET LINUX}
\newcommand{\Etudiant}{MAKEU TENKU STELY BELVA}
\newcommand{\Matricule}{24P826}
\newcommand{\Filiere}{Cybersécurité et Investigation Numérique}
\newcommand{\Niveau}{Niveau 4}
\newcommand{\Encadreur}{Ing. THIERRY MINKA}
\newcommand{\GradeEncadreur}{Enseignant -- ENSP}
\newcommand{\DateSoumission}{Mars 2026}

% ==============================================================================
\begin{document}
	
	% PAGE DE GARDE
	\input{PageGarde}
	
	% ==============================================================================
	% PAGE 1
	% ==============================================================================
	\thispagestyle{plain}
	\vspace*{0.2cm}
	
	% En-tête du résumé
	\begin{center}
		\begin{tikzpicture}
			\fill[CP!10, rounded corners=6pt] (0,0) rectangle (\textwidth, 1.3cm);
			\fill[CP, rounded corners=4pt] (0,0) rectangle (0.5cm, 1.3cm);
			\fill[CV, rounded corners=3pt] (0.5cm,0) rectangle (0.7cm, 1.3cm);
			\node[anchor=west, CB, font=\large\sffamily\bfseries]
			at (0.9cm, 0.65cm)
			{Exercices Fondamentaux d'Investigation Numérique};
			\node[anchor=east, CT!60, font=\small\sffamily\itshape]
			at (\textwidth - 0.3cm, 0.65cm);
		\end{tikzpicture}
	\end{center}
	
	\vspace{0.35cm}
	
	% ---------------------------------------------------------------
	% INTRODUCTION
	% ---------------------------------------------------------------
	\noindent
	Les exercices présentés ci-après constituent une mise en application
	directe des principes du \termecle{Processus Investigatif Universel} (PIU).
	Ils couvrent trois domaines opérationnels essentiels : l'identification
	des processus suspects sur un système live, l'investigation des mécanismes
	de persistance sous Windows, et l'analyse forensique de la
	\termecle{Master File Table} (MFT). Chaque exercice mobilise des outils
	concrets et des commandes éprouvées, dans le respect de la chaîne de
	traçabilité probatoire.
	
	\vspace{0.4cm}
	
	% ---------------------------------------------------------------
	% SECTION 1 : IDENTIFICATION DES PROCESSUS SUSPECTS
	% ---------------------------------------------------------------
	{\color{CP}\large\sffamily\bfseries\faAngleDoubleRight\enspace
		1. Identification des Processus Suspects sur Système Live}
	\vspace{0.15cm}
	\color{CV}\rule{\linewidth}{0.8pt}
	\color{CT}
	
	\vspace{0.3cm}
	
	\noindent
	L'identification des processus suspects constitue l'une des premières
	actions à mener lors d'une investigation sur un système en cours
	d'exécution (\termecle{live forensics}). Elle consiste à lister
	l'ensemble des processus actifs, à filtrer les entrées légitimes du
	noyau, puis à isoler ceux présentant des noms inhabituels ou des
	chemins d'exécution anormaux -- notamment des chemins situés dans
	des répertoires temporaires tels que \texttt{/tmp}, \texttt{/dev/shm}
	sous Linux, ou \texttt{\%APPDATA\%} sous Windows.
	
	\vspace{0.25cm}
	
	\textbf{Sur Linux :}
	
	\begin{lstlisting}[language=bash, caption={Listage et filtrage des processus suspects sous Linux}]
		# Lister les processus en excluant les processus noyau
		ps aux | grep -v '\[' | awk 'NR>1 {print $11}' | sort -u
		
		# Points d'analyse :
		# - Processus avec noms inhabituels
		# - Chemins d'execution suspects (/tmp, /dev/shm, /var/tmp)
		# - Processus sans chemin associe
	\end{lstlisting}
	
	\vspace{0.2cm}
	
	
	% IMAGE : Capture de la sortie de la commande ps aux filtrée
	%         montrant des processus suspects sous Linux
	\begin{figure}[h]
		     \centering
		     \includegraphics[width=0.9\textwidth]{ps_aux_linux}
		     \caption{Exemple de sortie filtrée de \texttt{ps aux} sous Linux}
		     \label{fig:ps_aux_linux}
		 \end{figure}
	% ---------------------------------------------------------------
	
	\vspace{0.2cm}
	
	\textbf{Sur Windows (PowerShell) :}
	
	\begin{lstlisting}[language=PowerShell, caption={Listage des processus hors répertoire système sous Windows}]
		# Lister les processus dont le chemin n'est pas dans C:\Windows
		Get-Process | Select-Object Name, Path, Id |
		Where-Object {$_.Path -notlike "C:\Windows\*"} |
		Sort-Object Name
	\end{lstlisting}
	
	\vspace{0.2cm}
	
	% ---------------------------------------------------------------
	% IMAGE : Capture PowerShell montrant la liste des processus
	%         filtrés hors C:\Windows\
	 \begin{figure}[h]
		    \centering
		     \includegraphics[width=0.9\textwidth]{get_process_windows}
		     \caption{Résultat de \texttt{Get-Process} filtré sous PowerShell}
		     \label{fig:get_process_windows}
		 \end{figure}
	% ---------------------------------------------------------------
	
	\vspace{0.35cm}
	
	% ---------------------------------------------------------------
	% SECTION 2 : PERSISTENCE SUR WINDOWS
	% ---------------------------------------------------------------
	{\color{CP}\large\sffamily\bfseries\faAngleDoubleRight\enspace
		2. Investigation des Mécanismes de Persistance sur Windows}
	\vspace{0.15cm}
	\color{CV}\rule{\linewidth}{0.8pt}
	\color{CT}
	
	\vspace{0.3cm}
	
	\noindent
	L'analyse des points de persistance permet de détecter toute tentative
	d'un programme malveillant de survivre à un redémarrage du système.
	Les vecteurs de persistance les plus courants sous Windows incluent
	les clés de registre \texttt{Run}, les tâches planifiées et les
	services configurés en démarrage automatique. Leur énumération
	systématique constitue une étape incontournable de la phase
	d'investigation analytique (E7--E10) du PIU.
	
	\vspace{0.25cm}
	
	\begin{lstlisting}[language=PowerShell, caption={Énumération des points de persistance communs sous Windows}]
		# Clés de registre Run -- niveau machine (tous les utilisateurs)
		Get-ItemProperty "HKLM:\SOFTWARE\Microsoft\Windows\CurrentVersion\Run"
		
		# Clés de registre Run -- niveau utilisateur courant
		Get-ItemProperty "HKCU:\SOFTWARE\Microsoft\Windows\CurrentVersion\Run"
		
		# Taches planifiees actives
		Get-ScheduledTask | Where-Object {$_.State -eq "Ready"}
		
		# Services configurés en démarrage automatique
		Get-Service | Where-Object {$_.StartType -eq "Automatic"}
	\end{lstlisting}
	
	\vspace{0.2cm}
	
	% ---------------------------------------------------------------
	% IMAGE : Capture PowerShell des clés de registre Run
	%         et des tâches planifiées suspectes détectées
	 \begin{figure}[h]
		     \centering
		     \includegraphics[width=0.9\textwidth]{persistence_registry}
		     \caption{Clés de registre \texttt{Run} et tâches planifiées analysées}
		     \label{fig:persistence_registry}
		 \end{figure}
		 
		 \begin{figure}[h]
		 	\centering
		 	\includegraphics[width=0.9\textwidth]{persistence_registry}
		 	\caption{Clés de registre \texttt{Run} et tâches planifiées analysées}
		 	\label{fig:persistence_registry_2}
		 \end{figure}
		 
		 \begin{figure}[h]
		 	\centering
		 	\includegraphics[width=0.9\textwidth]{persistence_registry}
		 	\caption{Clés de registre \texttt{Run} et tâches planifiées analysées}
		 	\label{fig:persistence_registry_3}
		 \end{figure}
		 
		 \begin{figure}[h]
		 	\centering
		 	\includegraphics[width=0.9\textwidth]{persistence_registry}
		 	\caption{Clés de registre \texttt{Run} et tâches planifiées analysées}
		 	\label{fig:persistence_registry_4}
		 \end{figure}
	% ---------------------------------------------------------------
	
	\vspace{0.35cm}
	
	% ---------------------------------------------------------------
	% SECTION 3 : ANALYSE DE LA MFT
	% ---------------------------------------------------------------
	{\color{CP}\large\sffamily\bfseries\faAngleDoubleRight\enspace
		3. Extraction et Analyse de la Master File Table (MFT)}
	\vspace{0.15cm}
	\color{CV}\rule{\linewidth}{0.8pt}
	\color{CT}
	
	\vspace{0.3cm}
	
	\noindent
	La \termecle{Master File Table} (MFT) est la structure centrale du
	système de fichiers NTFS. Elle contient une entrée pour chaque fichier
	et répertoire du volume, incluant les métadonnées temporelles issues
	de deux attributs distincts : \texttt{\$STANDARD\_INFORMATION} (0x10)
	et \texttt{\$FILE\_NAME} (0x30). La divergence entre ces deux horodatages
	constitue l'indicateur principal du \termecle{timestomping} -- technique
	de manipulation des dates à des fins de dissimulation forensique.
	
	\vspace{0.25cm}
	
	\noindent
	L'extraction s'effectue via \textbf{FTK Imager} (interface graphique)
	ou directement en ligne de commande avec \textbf{MFTECmd} d'Eric
	Zimmerman. Cette opération requiert des \textbf{privilèges Administrateur}
	pour accéder au fichier système \texttt{\$MFT} :
	
	\vspace{0.2cm}
	
	\begin{lstlisting}[language=bash, caption={Extraction de la MFT avec MFTECmd (privilèges admin requis)}]
		MFTECmd.exe -f "C:\$MFT" --csv "C:\output" --csvf mft.csv
	\end{lstlisting}
	
	\vspace{0.2cm}
	
	% ---------------------------------------------------------------
	% IMAGE : Capture du terminal montrant l'exécution réussie
	%         de MFTECmd avec la sortie CSV générée
	 \begin{figure}[h]
		     \centering
		     \includegraphics[width=0.9\textwidth]{mftecmd_output}
		     \caption{Exécution de \texttt{MFTECmd.exe} et génération du fichier CSV}
		     \label{fig:mftecmd_output}
		 \end{figure}
	% ---------------------------------------------------------------
	
	\vspace{0.25cm}
	
	\noindent
	Une fois le fichier CSV généré, le script Python suivant permet
	de filtrer et d'analyser automatiquement les artefacts forensiques
	d'intérêt : fichiers avec timestomping, fichiers supprimés, fichiers
	copiés et fichiers téléchargés depuis Internet.
	
	\vspace{0.2cm}
	
	\begin{lstlisting}[language=Python, caption={Analyse forensique du CSV MFT avec Python et pandas}]
		import pandas as pd
		
		df = pd.read_csv(r"C:\output\mft.csv")
		
		# ----------------------------------------------------------
		# 1. Détection du timestomping
		#    SI<FN = True indique une anomalie temporelle suspecte :
		#    $STANDARD_INFORMATION antérieur à $FILE_NAME
		# ----------------------------------------------------------
		stomped = df[df["SI<FN"] == True]
		print(f"[!] Fichiers avec timestomping détecté : {len(stomped)}")
		print(stomped[["FileName", "ParentPath",
		"LastModified0x10",
		"LastModified0x30"]].to_string())
		
		# ----------------------------------------------------------
		# 2. Fichiers supprimés (entrées MFT inactives)
		# ----------------------------------------------------------
		deleted = df[df["InUse"] == False]
		print(f"\n[!] Fichiers supprimés : {len(deleted)}")
		print(deleted[["FileName", "ParentPath",
		"LastModified0x10"]].to_string())
		
		# ----------------------------------------------------------
		# 3. Fichiers copiés (divergence Created0x10 / Created0x30)
		# ----------------------------------------------------------
		copied = df[df["Copied"] == True]
		print(f"\n[!] Fichiers copiés : {len(copied)}")
		print(copied[["FileName", "ParentPath",
		"Created0x10",
		"Created0x30"]].to_string())
		
		# ----------------------------------------------------------
		# 4. Fichiers téléchargés (présence d'un Zone.Identifier)
		# ----------------------------------------------------------
		downloaded = df[df["ZoneIdContents"].notna()]
		print(f"\n[!] Fichiers téléchargés depuis Internet : {len(downloaded)}")
		print(downloaded[["FileName", "ParentPath",
		"ZoneIdContents"]].to_string())
	\end{lstlisting}
	
	\vspace{0.2cm}
	
	% ---------------------------------------------------------------
	% IMAGE : Capture de la sortie Python montrant les fichiers
	%         suspects détectés (timestomping, supprimés, copiés,
	%         téléchargés)
	 \begin{figure}[h]
		     \centering
		     \includegraphics[width=0.9\textwidth]{python_mft_output}
		     \caption{Résultats de l'analyse forensique du CSV MFT par le script Python}
		     \label{fig:python_mft_output}
		 \end{figure}
	% ---------------------------------------------------------------
	
	\vspace{0.35cm}
	
	% ---------------------------------------------------------------
	% CONCLUSION
	% ---------------------------------------------------------------
	\begin{boxobjectifs}
		\textbf{Points clés à retenir :}
		\begin{itemize}
			\item Toujours travailler sur une \textbf{copie forensique} --
			jamais sur le système ou disque original
			\item Les droits \textbf{Administrateur} sont indispensables
			pour accéder aux fichiers système protégés (\texttt{\$MFT})
			\item La colonne \texttt{SI<FN} de MFTECmd est l'indicateur
			direct du \termecle{timestomping} -- pas besoin de le
			recalculer manuellement
			\item Les itérations entre phases d'acquisition et d'analyse
			sont normales et souhaitables dans le cadre du PIU
		\end{itemize}
	\end{boxobjectifs}
	
\end{document}